\lecture{23}{15. December 2025}{Exam 2023}

\begin{problem}{23.1}
  Show that
  \begin{equation*}
    f(x) = \begin{cases}
      1-x,  & 0 < x < 1\\
    -1-x, & -1 < x < 0
    \end{cases}
  \end{equation*}
  is an odd function.
\end{problem}

\begin{derivation}
  To determine if the function is odd or oven, we try inserting $-x$ as the argument to $f$.
  For $0<x<1$ we get:
  \begin{equation*}
    f(-x) = 1 + x = - f(x)
  .\end{equation*}
  And for $-1 < x < 0$ we get
  \begin{equation*}
    f(-x) = - 1 + x = -f(x)
  .\end{equation*}
  Hence the function is odd.
\end{derivation}

\newpage

\begin{problem}{23.2}
  Let $O_1(x)$ and $O_2(x)$ be odd functions and let $E(x)$ be an even function.
  Show that
  \begin{equation*}
    f(x) = O_1(x) \left( e^{\left| O_2 (x) E(x) \right|} + 1 \right)
  \end{equation*}
  is an odd function.
\end{problem}

\begin{derivation}
  We insert $-x$ into the function and get
  \begin{align*}
    f(-x) &= O_1(-x) \left( e^{\left| O_2(-x) E(-x) \right|} + 1 \right) \\
    &= - O_1(x) \left( e^{\left| - O_2(x) \cdot E(x) \right|} + 1 \right) \\
    &= - \left( O_1(x) \left( e^{\left| O_2(x) E(x) \right|} + 1 \right) \right) \\
    &= - f(x)
  .\end{align*}
  Hence, the function is odd.
\end{derivation}

\newpage

\begin{problem}{23.3}
  Find a period and the corresponding Fourier coefficients of
  \begin{enumerate}
    \item
      \begin{equation*}
        f(x) = \sin \left( \frac{\pi}{7} \right) + 2 \cos \left( \frac{2x}{\sqrt{2}\pi} \right) + 7 \sin \left( \frac{\sqrt{2}x}{3\pi} \right), \qquad - \infty < x < \infty
      .\end{equation*}
    \item
      \begin{equation*}
        f(x) = \sin \left( \sqrt{2}x \right) + \sin \left( 3 \sqrt{2}x \right), \qquad -\infty < x < \infty
      .\end{equation*}
  \end{enumerate}
\end{problem}

\begin{derivation}
  The general representation of a Fourier series is:
  \begin{equation*}
    f(x) = a_0 + \sum_{n = 1}^{\infty}  \left( a_n \cos \left( \frac{n \pi x}{L} \right) + b_n \sin \left( \frac{n \pi x}{L} \right) \right)
  .\end{equation*}
  
  \begin{enumerate}
    \item We rewrite our given function as
      \begin{equation*}
        f(x) = \underbrace{\sin \left( \frac{\pi}{7} \right)}_{a_0} + \underbrace{2}_{a_6} \cos \bigg( \frac{6 \pi x}{\underbrace{3 \sqrt{2}\pi^2}_L} \bigg) + \underbrace{7}_{b_2} \sin \bigg( \frac{2 \pi x}{\underbrace{3\sqrt{2} \pi^2}_L} \bigg)
      .\end{equation*}
      Hence, $p = 2 L = 6 \sqrt{2} \pi^2$, $a_0 = \sin (\pi / 7)$, $a_6 = 2$, $b_2 = 7$ and all other Fourier coefficients are zero.
    \item Once again, we rewrite as
      \begin{equation*}
        f(x) = \underbrace{1}_{b_2} \cdot \sin \bigg( \frac{2\pi x}{ \underbrace{\sqrt{2}\pi}_{L}} \bigg) + \underbrace{1}_{b_6} \sin \bigg( \frac{6 \pi x}{\underbrace{\sqrt{2}\pi}_{L}} \bigg) 
      .\end{equation*}
      Hence, $p = 2L = 2 \sqrt{2} \pi$, $b_2 = 1$, $b_6 = 1$ and all other Fourier coefficients are zero.
  \end{enumerate}
\end{derivation}

\finalresult{
\begin{aligned}
  &\mathrm{a}) & p &= 6 \sqrt{2}\pi^2, \quad a_0 = \sin \frac{\pi}{7}, \quad a_6 = 2, \quad b_2 = 7 \\
  & \mathrm{b}) & p &= 2 \sqrt{2}\pi, \quad b_2 = 1,\quad b_6 = 1
.\end{aligned}
}

\begin{problem}{23.4}
  Consider the function
  \begin{equation*}
    f(x) = \sin(2x), 0 \leq x \leq \pi
  .\end{equation*}
  
  \begin{enumerate}
    \item Find the period and the corresponding Fourier coefficients of the expansion of $f(x)$.
    \item Find the period and the corresponding Fourier coefficients of the odd expansion of $f(x)$.
  \end{enumerate}
\end{problem}

\begin{derivation}
  \begin{enumerate}
    \item As $\sin$ is in odd function, the expansion of $f(x) = \sin (2x), 0 \leq x \leq \pi$ is simply $f_0(x) = \sin(2x), - \pi / 2 \leq x \leq \pi / 2$.
      Rewriting this we get
      \begin{equation*}
        f_0(x) = \underbrace{1}_{b_1} \cdot \sin \frac{\pi x}{\underbrace{\frac{\pi}{2}}_L}
      .\end{equation*}
      Meaning $p = 2L = \pi$ and $b_1 = 1$ and all other Fourier coefficients are zero.
    \item The odd expansion of an odd function will also yield the same odd function, however, for the odd expansion the period will be twice as long, as $f_2 (x) = \sin(2x), - \pi \leq x \leq \pi$.
      Rewriting, we get
      \begin{equation*}
        f_2(x) = \underbrace{1}_{b_2} \cdot \sin \frac{2\pi x}{\underbrace{\pi}_L}
      .\end{equation*}
      Hence, $p = 2L = 2\pi$, $b_2 = 1$ and all other Fourier coefficients are zero.
  \end{enumerate}
\end{derivation}

\finalresult{
\begin{aligned}
  &\mathrm{a}) & p &= \pi, \quad b_1 = 1 \\
  &\mathrm{b}) & p &= 2\pi, \quad b_2 = 1
.\end{aligned}
}

\begin{problem}{23.5}
  Rewrite the PDE
  \begin{equation*}
    u_{y y}(x,y) + u_x(x,y) = 0
  \end{equation*}
  in terms of the variables
  \begin{equation*}
    p(x,y) = e^{y}, q(x,y) = e^{x}
  .\end{equation*}
\end{problem}

\begin{derivation}
  We define
  \begin{equation*}
    \overline{u}(p,q) = \overline{u}(p(x,y) , q(x,y) ) = u(x,y) \implies \overline{u}_{yy}(p,q) + \overline{u}_x(p,q) = 0
  .\end{equation*}
  The first derivative with respect to $x$ is
  \begin{align*}
    \overline{u}_x(p,q) &= \overline{u}_p(p,q) p_x(x,y) + \overline{u}_q(p,q) q_x(x,y) \\
    &= \overline{u}_p(p,q) \cdot 0 + \overline{u}_q(p,q) \cdot e^{x} \\
    &= \overline{u}_q(p,q) \cdot q
  .\end{align*}
  The first derivative with respect to $y$ is
  \begin{align*}
    \overline{u}_y(p,q) &= \overline{u}_p(p,q) p_y(x,y) + \overline{u}_q(p,q) q_y(x,y) \\
    &= \overline{u}_p(p,q) \cdot e^{y} + \overline{u}_q(p,q) \cdot 0 \\
    &= \overline{u}_p(p,q) \cdot p
  .\end{align*}
  The second derivative with respect to $y$ is
  \begin{align*}
    \overline{u}_{yy}(p,q) &= \overline{u}_{pp}(p,q) p_y((x,y) \cdot e^{y} + \overline{u}_{p} (p,1) e^{y} \\
    &=  \overline{u}_{pp}(p,q) p^2 + \overline{u}_p(p,q) p
  .\end{align*}
  By insertion we get
\end{derivation}

\finalresult{
\overline{u}_{pp}(q,p) p^2 + \overline{u}_p(q,p) p + \overline{u}_q(p,q) q = 0
}

\begin{problem}{23.6}
  Consider the PDE
  \begin{equation*}
    u_{x x}(x,t) + u_{t t}(x,t) = u(x,t), 0 \leq x \leq \infty, t \geq 0
  \end{equation*}
  and let $\hat{u}_C(\omega, t) = \mathcal{F}_C (u(x,t))$ be the Fourier cosine transform of $u(x,t)$ with respect to $x$.
  Derive an ODE for $\hat{u}_C(\omega,t)$ with respect to $t$.
  You may assume that
  \begin{equation*}
    \mathcal{F}_C(u_{t t}(x,t)) = \frac{\partial ^2}{\partial t^2} \mathcal{F}_C \left( u(x,t) \right)
  \end{equation*}
  and $u_x(0,t) = 0$ for all $t$.
\end{problem}

\begin{derivation}
  We know that $\mathcal{F}_C(f'') = -\omega^2 \mathcal{F}_C(f) - \sqrt{2 / \pi} f'(0)$.
  Hence, we can rewrite our expression as
  \begin{align*}
    - \omega^2 \hat{u}_C(\omega, t) - \sqrt{\frac{2}{\pi}} \cdot 0 + \frac{\partial ^2}{\partial t^2} \hat{u}_C(\omega,t ) &= \hat{u}_C (\omega, t) \\
    - \omega^2 \hat{u}_C(\omega, t) + \frac{\partial ^2}{\partial t^2} \hat{u}_C(\omega, t) &= \hat{u}_C(\omega, t)
  .\end{align*}
\end{derivation}

\finalresult{
- \omega^2 \hat{u}_C(\omega, t) + \frac{\partial ^2}{\partial t^2} \hat{u}_C(\omega, t) = \hat{u}_C(\omega, t)
}

\begin{problem}{23.7}
  Consider the PDE
  \begin{equation*}
    u_{y y}(x,y) + u_x(x,y) = 0, 0 \leq y \leq 1, x \geq 0
  \end{equation*}
  subject to the conditions
  \begin{equation*}
    u(x,0) = u(x,1) = 0 \text{ for all } x
  \end{equation*}
  and
  \begin{equation*}
    u(0,y) = 3 \sin (2\pi y) + 2 \sin (3\pi y) \text{ for } 0 \leq y \leq 1
  .\end{equation*}
  
  Find the solution $u(x,y)$.
\end{problem}

\begin{derivation}
  Using the method of separation of variables, we assume $u(x,y) = F(x) Q(y)$ and get
  \begin{align*}
    F(x) Q''(y) + F'(x) Q(y) &= 0 \\
    \frac{Q''(y)}{Q(y)} &= - \frac{F'(x)}{F(x)} = k
  .\end{align*}
  Which leaves us with the two ODEs
  \begin{align*}
    Q''(y) - kQ(y) &= 0 \\
    F'(x) - k F(x) &= 0
  .\end{align*}
  Our boundary conditions gives us
  \begin{align*}
    u(x,0) &= F(x) Q(0) = 0 \\
    \implies Q(0) &= 0 \\
    u(x,1) &= F(x) Q(1) = 0 \\
    \implies Q(1) &= 0
  .\end{align*}

  Assuming $k = 0 \implies F(x) = Q(y) = 0$ which is an uninteresting case.
  
  If we assume $k > 0$ then
  \begin{gather*}
    Q(y) = A e^{\sqrt{k}y} + B e^{- \sqrt{k}y} \\
    Q(0) = 0 \implies A = - B \\
    Q(1) = 0 \implies 0 = A e^{\sqrt{k}} + Be^{-\sqrt{k}} \implies A = 0 = B
  .\end{gather*}
  This is also uninteresting.

  Hence, our only interesting case arises for $k < 0$, for which we get
  \begin{gather*}
    Q(y) = A \cos \left( \sqrt{-k} y \right) + B \sin \left( \sqrt{-k}y \right) \\
    Q(0) = 0 \implies A = 0 \\
    Q(1) = 0 \implies B \sin \left( \sqrt{-k} \right)  = 0
  .\end{gather*}
  We set $\sqrt{-k} = p$ and get
  \begin{equation*}
    Q(1) = B \sin p = 0 \implies p = p_n = n \pi, n = 1,2,\ldots 
  .\end{equation*}
  We choose $B = 1$ and get
  \begin{equation*}
    Q_n(y) = \sin (n \pi y)
  .\end{equation*}
  From the Lecture notes we have that
  \begin{equation*}
    F_n(x) = B_n e^{n^2 \pi^2 x}
  .\end{equation*}
  Hence,
  \begin{align*}
    u_n (x,y) &= F_n(x) Q_n(y) \\
    &= B_n e^{n^2 \pi^2 x} \sin (n \pi y) \\
    u(x,y) &= \sum_{n = 1}^{\infty} u_n(x,y) \\
    u(x,y) &= \sum_{n = 1}^{\infty} B_n e^{n^2 \pi^2 x} \sin (n \pi y)
  .\end{align*}
  We now compare this with our initial condition:
  \begin{gather*}
    u(0,y) = 3 \sin (2\pi y) 2 \sin (3\pi y) \\
    \implies B_2 = 3 \quad B_3 = 2
  \end{gather*}
  and all other Fourier coefficients are zero.

  Hence,
\end{derivation}

\finalresult{
 u(x,y) = 3 e^{4\pi^2 x} \sin (2\pi y) +  2 e^{9\pi^2 x} \sin (3\pi y)
}

\begin{problemwithimage}{e23_8}{23.8}
  Consider the two-dimensional Laplace equation
  \begin{equation*}
    u_{x x}(x,y) + u_{yy}(x,y) = 0, 0 \leq x \leq 2, 0 \leq y \leq 2
  \end{equation*}
  subject to the mixed boundary conditions
  \begin{align*}
    u_x(0,y) &= 0, & 0 \leq &y \leq 2 \\
    u_x(2,y) &= 0, & 0 \leq &y \leq 2 \\
    u(x,2) &= 1, & 0 \leq &x \leq 2 \\
    u(x,0) &= 0, & 0 \leq &x \leq 2
  .\end{align*}
  Derive the corresponding system of difference equations, using mesh size $h = 1$.
\end{problemwithimage}

\begin{derivation}
  We have the known boundary values
  \begin{align*}
    u_{00} &= 0, & u_{10} &= 0, & u_{20} &= 0, \\
    u_{02} &= 1, & u_{12} &= 1, & u_{22} &= 1
  .\end{align*}
  And a corresponding difference equation of
  \begin{equation*}
    u_{i+1,j} + u_{i,j+1} + u_{i-1,j} + u_{i,j-1} - 4u_{i,j} = 0
  .\end{equation*}
  
  We start by calculating the unknown meshpoint $P_{01}$ using the difference equation as
  \begin{align*}
    u_{1,1} + u_{0,2} + u_{-1,1} + u_{0,0} - 4u_{0,1} &= 0 \\
    u_{1,1} + 1 + u_{-1,1} + 0 - 4u_{0,1} &= 0 \\
    u_{1,1} +u_{-1,1} - 4u_{0,1} &= -1
  .\end{align*}
  We find $u_{-1,1}$ using the directional derivative as
  \begin{equation*}
    0 = \frac{u_{1,1} - u_{-1,1}}{2} \implies u_{-1,1} = u_{1,1}
  .\end{equation*}
  Hence,
  \begin{equation*}
    2u_{11} - 4u_{01} = -1
  .\end{equation*}
  
  For the meshpoint $P_{11}$ we get
  \begin{align*}
    u_{21} + u_{12} + u_{01} + u_{10} - 4u_{11} &= 0 \\
    u_{21} + 1 + u_{01} + 0 - 4u_{11} &= 0 \\
    u_{21} + u_{01} -4u_{11} &= -1
  .\end{align*}

  For the meshpoint $P_{21}$ we get
  \begin{align*}
    u_{31} + u_{22} + u_{11} + u_{20} - 4u_{21} &= 0 \\
    u_{31} + 1 + u_{11} + 0 - 4u_{21} &= 0 \\
    u_{31} + u_{11} - 4u_{21} &= -1
  .\end{align*}
  We find $u_{31}$ using the directional derivative as
  \begin{equation*}
    0 = \frac{u_{31} - u_{11}}{2} \implies u_{31} = u_{11}
  .\end{equation*}
  Hence
  \begin{equation*}
    2u_{11} - 4u_{21} = -1
  .\end{equation*}
\end{derivation}

\finalresult{
\begin{aligned}
  2u_{11} - 4u_{01} &= -1 \\
  u_{21} + u_{01} - 4u_{11} &= -1 \\
  u_{31} + u_{11} - 4u_{21} &= -1
.\end{aligned}
}
